\documentclass[11pt]{book}
\usepackage{fontspec}
\setmainfont{Latin Modern Roman}
\setsansfont{Latin Modern Sans}
\usepackage{amsmath,amssymb,mathtools}
\usepackage{unicode-math}
\setmathfont{Latin Modern Math}
\usepackage{xurl}
\usepackage[hidelinks]{hyperref}
\newcommand{\sref}[2]{\hyperlink{source:#1:#2}{[S#2]}}
\newcommand{\cataloguescope}[1]{\paragraph{Mathematical question.}#1}
\begin{document}
\section{Research Problems in Function Theory — Problem 4.15}\label{statement}

\subsection{Definitions and mathematical statement}

\paragraph{Shared notation.}$\mathbb N=\{1,2,\ldots\}$, and $\mathbb Z,\mathbb Q,\mathbb R,\mathbb C$ denote the integers, rationals, reals and complex numbers. Any other conventions are specified locally.


Write $\mathbb D=\{z\in\mathbb C:|z|<1\}$ and $\mathbb T=\partial\mathbb D$.For each $n\ge1$, consider all $2^n$ polynomials
\[P_\epsilon(z)=\sum_{k=1}^n\epsilon_kz^k,\qquad\epsilon=(\epsilon_1,\ldots,\epsilon_n)\in\{-1,1\}^n.
\]
Let $Z_{\mathbb D}(P_\epsilon)$ be the number of zeros in the open disc $\mathbb D$, counted with multiplicity and including the zero at the origin.
\paragraph{Statement recorded in the supplied source.}
Is it true that for every $\delta>0$,
\[2^{-n}\#\left\{\epsilon\in\{-1,1\}^n:\left|Z_{\mathbb D}(P_\epsilon)-\tfrac n2\right|>\delta n\right\}\longrightarrow0\quad(n\to\infty)?\]
Equivalently, all but $o(2^n)$ of these polynomials have $n/2+o(n)$ zeros in $\mathbb D$, with the error interpreted uniformly outside a vanishing fraction of sign choices. \sref{statement}{2}
\subsection{Short English statement}
For almost every sign choice, does a polynomial with coefficients plus or minus one and no constant term have about half its zeros inside the unit disc?
\subsection{Sources}
\begin{description}
\item[\hypertarget{source:statement:1}{S1}] ULAM.AI and the UnsolvedMath Contributors, \emph{UnsolvedMath: A Curated Collection of Open Mathematics Problems}, catalogue entry. CC BY 4.0. \url{https://huggingface.co/datasets/ulamai/UnsolvedMath}.
\item[\hypertarget{source:statement:2}{S2}] Walter K. Hayman and Edward F. Lingham, \emph{Research Problems in Function Theory}, arXiv:1809.07200, 2018. Problem 4.15, its complete Update 4.15, and the relevant chapter notation. \url{https://arxiv.org/pdf/1809.07200}.
\end{description}
\label{end:statement}
\end{document}
